% Generated draft: requires mathematical and source review.
\documentclass[11pt]{article}
\usepackage[margin=25mm]{geometry}
\usepackage{fontspec}
\setmainfont{Latin Modern Roman}
\usepackage{amsmath,amssymb,mathtools,unicode-math}
\setmathfont{Latin Modern Math}
\usepackage{enumitem,needspace,xurl,hyperref}
\setlength{\parindent}{0pt}
\setlength{\parskip}{5pt}
\setlength{\emergencystretch}{3em}
\newcommand{\N}{\mathbb{N}}
\newcommand{\Z}{\mathbb{Z}}
\newcommand{\Q}{\mathbb{Q}}
\newcommand{\R}{\mathbb{R}}
\newcommand{\C}{\mathbb{C}}
\newcommand{\F}{\mathbb{F}}
\newcommand{\A}{\mathbb{A}}
\newcommand{\PP}{\mathbb{P}}
\newcommand{\E}{\mathbb{E}}
\newcommand{\eps}{\varepsilon}
\newcommand{\dd}{\,\mathrm{d}}
\begin{document}
\section*{UnsolvedMath ID: 800334}
\textbf{Source problem code:} catalogue entry
\subsection{Definitions and mathematical statement}
For an integer matrix \(M\in\operatorname{Mat}_{m\times n}(\Z)\), define its cokernel by
\[
\operatorname{coker}_{\Z}(M)=\Z^m/M\Z^n.
\]
Here the matrix represents a homomorphism from \(\Z^n\) to \(\Z^m\). The integral interpretation is suggested by the source's reference to a non-cyclic group, but the coefficient ring is not explicitly specified in the supplied record.

Write \(C_{a,b,c}\) for the Carlitz matrix intended by the source and \(K_{a,b,c}\) for its intended Kasteleyn matrix associated with an \(a,b,c\) hexagon. These are placeholders for the source's matrix families: the supplied record gives neither their entries and dimensions nor the admissible parameter triples, the precise hexagonal region, or the Kasteleyn conventions.

\textbf{Problem (Propp, 1998, Problem 5).} Determine, up to isomorphism, the abelian groups
\[
\operatorname{coker}_{\Z}(C_{a,b,c})
\]
for all admissible parameter triples \((a,b,c)\). The source describes this as equivalent to determining \(\operatorname{coker}_{\Z}(K_{a,b,c})\). If possible, give an interpretation of the cokernel in terms of the tilings of the corresponding hexagon.

The supplied source additionally reports that the cokernel for \(a=b=c=2\) is non-cyclic. This is retained as a source assertion, without an independently verified group identification. The claimed equivalence of the two matrix formulations is likewise retained as supplied; the record does not specify the correspondence establishing it. A fully explicit formulation requires the missing matrix, parameter, region, and tiling definitions.
\subsection{Short English statement}
Determine the cokernels of the Carlitz matrices, described by Propp as equivalently the cokernels of Kasteleyn matrices for \(a,b,c\) hexagons, and, if possible, interpret these groups through the tilings. The source reports a non-cyclic cokernel when \(a=b=c=2\). The supplied record lacks the definitions needed to specify these matrix families completely.
\subsection{Sources}
\begin{itemize}
\item Local UnsolvedMath registry: supplied complete source record, permanent ID \(800334\), exact source problem code \texttt{catalogue entry}, title \emph{Enumeration of Matchings --- Problem 5}. The registry attributes the problem to Propp, with proposed year 1998, and labels it \texttt{open}; this is a supplied status label, not an independently verified assessment.
\item Propp, \emph{Twenty Open Problems in Enumeration of Matchings} (1998), Problem 5 (\#5). Supplied source URL: \url{https://arxiv.org/abs/math/9801060}. The provenance records access on 2026-10-16 and extraction by \texttt{arxiv\_latex}; these are supplied provenance statements, not a claim of independent source access or verification.
\end{itemize}
\end{document}
